% year: תשע"ח
% type: מבחן
% semester: ב
% moed: א
% number: 3א
% points: 15
% categories: definite-integrals
% source: exams/מבחנים/תשעח/מועד א תשעח סמסטר ב.pdf

\begin{question}
חשבו: $\displaystyle\int_{-1}^15x^4\big(\ln(x+2)+\sin x\big)dx$
\end{question}

\begin{hint}
$5x^4\sin x$ היא פונקציה אי-זוגית, ולכן האינטגרל שלה על $[-1,1]$ מתאפס. לחלק השני, אינטגרציה בחלקים עם $u=\ln(x+2)$, $dv=5x^4dx$.
\end{hint}

\begin{solution}
\textbf{החלק עם הסינוס.} $g(x)=5x^4\sin x$ אי-זוגית ($g(-x)=-g(x)$) ורציפה, ולכן $\int_{-1}^1 5x^4\sin x\,dx=0$.

\textbf{החלק עם הלוגריתם.} אינטגרציה בחלקים עם $u=\ln(x+2)$, $dv=5x^4dx$, $du=\frac{dx}{x+2}$, $v=x^5$:
\[
\int_{-1}^15x^4\ln(x+2)\,dx=\Big[x^5\ln(x+2)\Big]_{-1}^1-\int_{-1}^1\frac{x^5}{x+2}dx=\ln3-\int_{-1}^1\frac{x^5}{x+2}dx
\]
(כי ב-$x=-1$: $(-1)^5\ln1=0$). חילוק פולינומים:
\[
\frac{x^5}{x+2}=x^4-2x^3+4x^2-8x+16-\frac{32}{x+2}.
\]
על הקטע הסימטרי $[-1,1]$ האינטגרלים של החזקות האי-זוגיות מתאפסים, ולכן
\[
\int_{-1}^1\frac{x^5}{x+2}dx=\int_{-1}^1\left(x^4+4x^2+16\right)dx-32\int_{-1}^1\frac{dx}{x+2}
=\frac25+\frac83+32-32\ln3=\frac{526}{15}-32\ln3.
\]
\textbf{תשובה:}
\[
\int_{-1}^15x^4\big(\ln(x+2)+\sin x\big)dx=\ln3-\frac{526}{15}+32\ln3=33\ln3-\frac{526}{15}\approx1.1875.
\]
\end{solution}
